\chapter{Frontend Architecture}
\label{ch:frontend-architecture}

The \sysname{} frontend is a Next.js 16 application using the App Router,
React 19, Tailwind CSS v4, and shadcn/ui.  This chapter covers the routing
structure, the Nebula design system, component patterns, graph visualization,
and state management.


% ─────────────────────────────────────────────────────────────────────
\section{Next.js 16 App Router}
\label{sec:app-router}

\sysname{} uses the App Router (the \ic{app/} directory), not the legacy
Pages Router.  Key concepts:

\begin{itemize}
    \item \textbf{Layouts} (\ic{layout.tsx}): Wrap child routes with shared
          UI (sidebar, navigation, providers).  The root layout at
          \ic{app/layout.tsx} applies the Instrument Sans font, theme
          provider, session provider, and React Query provider.

    \item \textbf{Pages} (\ic{page.tsx}): The actual content for a route.
          Each route segment maps to a directory, and the \ic{page.tsx}
          inside it renders the page content.

    \item \textbf{Route groups} (\ic{(group)/}): Organize routes without
          affecting the URL.  For example, \ic{app/(dashboard)/} groups
          all authenticated pages under a shared layout without adding
          ``dashboard'' to the URL path.

    \item \textbf{Server Components}: By default, every component in the
          \ic{app/} directory is a React Server Component.  Components
          that use browser APIs, hooks, or event handlers must be marked
          with \ic{'use client'} at the top of the file.
\end{itemize}

\paragraph{Adding a new page.}
To add a page at \ic{/settings}:

\begin{enumerate}
    \item Create \ic{app/(dashboard)/settings/page.tsx}.
    \item Export a default React component.
    \item The route inherits the dashboard layout automatically.
    \item Add a navigation link in \ic{components/layout/Sidebar.tsx}.
\end{enumerate}


% ─────────────────────────────────────────────────────────────────────
\section{Nebula Design System}
\label{sec:nebula-design-system}

\sysname{}'s visual identity is called ``Nebula''---a dark-first theme
built on violet, rose, and orange accent colors with glassmorphism effects.
All colors are defined as CSS custom properties in \codefile{app/globals.css}
and consumed via Tailwind's semantic classes.

\subsection{Semantic Color Tokens}

\begin{table}[ht]
\centering
\caption{Semantic color tokens (light / dark).}
\label{tab:semantic-colors}
\begin{tabularx}{\textwidth}{l l l X}
\toprule
\textbf{Token} & \textbf{Light} & \textbf{Dark} & \textbf{Usage} \\
\midrule
\ic{--primary}      & Violet 48\%  & Violet 58\%  & Primary actions, focus rings \\
\ic{--accent}       & Rose 50\%    & Rose 60\%    & Secondary accents, highlights \\
\ic{--destructive}  & Red 51\%     & Red 45\%     & Delete, error states \\
\ic{--muted}        & Gray 94\%    & Dark 14\%    & Disabled, secondary text \\
\ic{--background}   & Off-white    & Space 6\%    & Page background \\
\ic{--card}         & White        & Glass 10\%   & Card surfaces \\
\bottomrule
\end{tabularx}
\end{table}

\subsection{Entity Type Colors}

Entities in the graph and in badges are color-coded by type:

\begin{table}[ht]
\centering
\caption{Entity type color assignments.}
\label{tab:entity-colors}
\begin{tabular}{lll}
\toprule
\textbf{Entity Type} & \textbf{Color} & \textbf{Hex} \\
\midrule
Technology & Orange  & \ic{\#fb923c} \\
Pattern    & Rose    & \ic{\#ec4899} \\
Concept    & Violet  & \ic{\#a78bfa} \\
Person     & Emerald & \ic{\#34d399} \\
System     & Green   & \ic{\#4ade80} \\
\bottomrule
\end{tabular}
\end{table}

These colors are used consistently across the graph visualization,
dashboard charts, entity badges, and search result cards.


% ─────────────────────────────────────────────────────────────────────
\section{Component Patterns}
\label{sec:component-patterns}

The codebase follows several conventions that ensure consistency:

\subsection{\texttt{cn()} for Class Names}

All className composition uses the \ic{cn()} utility from
\codefile{lib/utils.ts}, which merges Tailwind classes with
\ic{clsx} and resolves conflicts with \ic{tailwind-merge}:

\begin{lstlisting}[style=typescript, caption={Using \ic{cn()} for conditional classes.}]
<div className={cn(
  "rounded-lg border p-4",
  variant === "glass" && "backdrop-blur-xl bg-card/50",
  className
)} />
\end{lstlisting}

Never concatenate class strings manually---\ic{cn()} handles
deduplication and conflict resolution.

\subsection{CVA for Variants}

Component variants are defined using Class Variance Authority (CVA):

\begin{lstlisting}[style=typescript, caption={CVA variant definition for a Button.}]
const buttonVariants = cva("base-classes rounded-lg font-medium", {
  variants: {
    variant: {
      default: "bg-primary text-primary-foreground",
      glass: "backdrop-blur-xl bg-card/50 border-white/10",
      gradient: "bg-gradient-to-r from-violet-500 via-fuchsia-500 ...",
      destructive: "bg-destructive text-destructive-foreground",
    },
    size: {
      default: "h-10 px-4",
      sm: "h-8 px-3 text-sm",
      lg: "h-12 px-6 text-lg",
    },
  },
  defaultVariants: { variant: "default", size: "default" },
})
\end{lstlisting}

\subsection{\texttt{React.forwardRef}}

All UI components use \ic{React.forwardRef} so that parent components
can attach refs for focus management, scroll control, and animations:

\begin{lstlisting}[style=typescript, caption={ForwardRef pattern.}]
const Card = React.forwardRef<HTMLDivElement, CardProps>(
  ({ className, variant, ...props }, ref) => (
    <div ref={ref} className={cn(cardVariants({ variant }), className)}
         {...props} />
  )
)
Card.displayName = "Card"
\end{lstlisting}

\subsection{Compound Components}

Complex components export multiple related sub-components:
\ic{Card}, \ic{CardHeader}, \ic{CardTitle}, \ic{CardContent},
\ic{CardFooter}.  This gives consumers full control over layout without
exposing internal implementation details.


% ─────────────────────────────────────────────────────────────────────
\section{Graph Visualization}
\label{sec:graph-visualization}

The knowledge graph is rendered using React Flow (\ic{@xyflow/react})
with custom node types and automatic layout via Dagre.

\begin{itemize}
    \item \textbf{Custom node types}:
          \begin{itemize}
              \item \emph{Decision nodes} render as gradient cards with a
                    violet-to-rose background, showing the trigger text and
                    a confidence badge.
              \item \emph{Entity nodes} render as colored pills, tinted
                    by entity type (orange for technology, rose for pattern,
                    etc.).
          \end{itemize}

    \item \textbf{Layout algorithms}: Dagre provides four layout modes:
          force-directed, hierarchical, clustered, and radial.  The user
          can switch between them from the graph toolbar.

    \item \textbf{Interaction}: Arrow keys for panning, Tab/Enter for
          node selection, Escape to deselect.  A minimap and zoom
          controls are always visible.

    \item \textbf{Performance}: Large graphs (500+ nodes) use
          \ic{React.memo} with custom comparison functions on node
          components to prevent unnecessary re-renders.
\end{itemize}


% ─────────────────────────────────────────────────────────────────────
\section{State Management}
\label{sec:state-management}

\sysname{} does not use a global state library like Redux.  Instead,
it combines several purpose-specific tools:

\begin{table}[ht]
\centering
\caption{State management strategy.}
\label{tab:state-management}
\begin{tabularx}{\textwidth}{l l X}
\toprule
\textbf{Kind of State} & \textbf{Tool} & \textbf{Details} \\
\midrule
Server / async &
    React Query &
    60-second stale time.  Queries auto-refetch on window focus.
    Mutations invalidate related queries. \\
\addlinespace
Form &
    react-hook-form + Zod &
    Schema-first validation.  Zod schemas are shared with
    backend Pydantic models where possible. \\
\addlinespace
Theme &
    next-themes &
    System / light / dark.  The \ic{.dark} class on
    \ic{<html>} toggles CSS custom properties. \\
\addlinespace
Auth session &
    NextAuth v5 &
    \ic{SessionProvider} at the root layout.
    \ic{useSession()} hook in client components. \\
\addlinespace
Local / ephemeral &
    \ic{useState} / \ic{useReducer} &
    Component-level state for UI toggles, form drafts,
    modal open/close. \\
\bottomrule
\end{tabularx}
\end{table}


% ─────────────────────────────────────────────────────────────────────
\section{Gotchas}
\label{sec:frontend-gotchas}

Several non-obvious pitfalls have tripped up contributors.  These are
documented here to save you debugging time.

\begin{warning}
\textbf{Radix ScrollArea breaks \texttt{scrollTop}.} Radix's
\ic{ScrollArea} component sets \ic{overflow: hidden} on its root element,
which prevents programmatic scroll control via \ic{scrollTop}.  If you
need auto-scrolling (e.g., the chat panel in the Ask view), use a plain
\ic{<div>} with \ic{overflow-y-auto} instead.
\end{warning}

\begin{warning}
\textbf{React Flow hook ordering.}  The \ic{useNodesState} and
\ic{useEdgesState} setter functions must be declared \emph{before} any
\ic{useCallback} that references them.  React hooks must always be called
in the same order, and referencing a setter before its declaration
causes stale-closure bugs that are silent in development but break in
production.
\end{warning}

\begin{warning}
\textbf{CSS \texttt{:root} in dark mode.}  Because \ic{next-themes}
applies the \ic{.dark} class on \ic{<html>}, bare \ic{:root} selectors
apply in \emph{both} light and dark mode.  Use \ic{:root:not(.dark)} for
light-mode-only overrides.  Getting this wrong causes colors to bleed
between themes.
\end{warning}

\begin{tryit}
Explore the design system live: start the frontend (\ic{pnpm dev:web}),
navigate to the graph view, and try switching between layout modes.
Open the browser DevTools and inspect a decision node to see the Nebula
color tokens in action.
\end{tryit}
